% GENERATED FILE. Edit canonical lessons, then run npm run generate:books.
% canonical-source-hash: fnv1a64:9a4488dc7572bd97
% canonical-lessons: JA-C205-madoguchi, JA-C205-ekiin, JA-C205-kaisatsuguchi, JA-C205-katamichi, JA-C205-oufuku

\chapter{At the Ticket Window}
\label{ch:ja-at-the-ticket-window}
\hlchaptermodality{205}

\begin{chapteropening}
\textbf{By the end of this chapter:} \emph{I can buy a ticket and find the ticket gate.}

Complete the last lesson of chapter 205: I can buy a ticket and find the ticket gate.
\end{chapteropening}

\section[\texorpdfstring{\ja{まどぐち} (madoguchi)}{まどぐち (madoguchi)}]{\ja{まどぐち} (madoguchi) — a ticket window; a service counter}

\label{lesson:JA-C205-madoguchi}



\begin{quote}
Before the new one: say the Japanese for to drop in, then the Japanese for lodging.
\end{quote}

\subsection*{You'll want to know: \ja{まどぐち}}
\textbf{\ja{まどぐち}} — \emph{madoguchi} — \textquotedblleft{}a ticket window; a service counter\textquotedblright{}.

\begin{quote}
\textbf{\ja{まどぐちで} \ja{きっぷを} \ja{かいます。}} — \emph{madoguchi de kippu o kaimasu} — \textquotedblleft{}I buy a ticket at the ticket window.\textquotedblright{}
\end{quote}

\begin{cousinweb}
\textbf{\ja{まどぐち}} (\emph{madoguchi}) is native Japanese (wago), built from \emph{mado}, \textquotedblleft{}window\textquotedblright{}, and \emph{kuchi}, \textquotedblleft{}mouth, opening\textquotedblright{}, both inherited rather than borrowed; \emph{kuchi} is voiced to \emph{guchi} (JMdict, KANJIDIC2).
\end{cousinweb}

\begin{grammarlens}[title={What you've built}]
One more word for everyday life at A2.
\end{grammarlens}

\subsection*{Your turn}
\begin{itemize}
\raggedright
  \item \emph{Say it:} \emph{madoguchi}
  \item \emph{Say it:} \emph{madoguchi}, once more
  \item \emph{Say it:} \emph{madoguchi de kippu o kaimasu} — I buy a ticket at the ticket window.
  \item \emph{Recall:} say the Japanese for to drop in, then the Japanese for lodging, then say \emph{madoguchi} again
\end{itemize}

\begin{tcolorbox}[breakable,colback=teal!4,colframe=teal!35!black,title={Before you move on}]
What does \ja{まどぐち} mean? (\textquotedblleft{}A ticket window; a service counter\textquotedblright{}.) Say it once more.
\end{tcolorbox}
\section[\texorpdfstring{\ja{えきいん} (ekiin)}{えきいん (ekiin)}]{\ja{えきいん} (ekiin) — a station attendant}

\label{lesson:JA-C205-ekiin}



\begin{quote}
Before the new one: say the Japanese for lodging, then the Japanese for a ticket window.
\end{quote}

\subsection*{You'll want to know: \ja{えきいん}}
\textbf{\ja{えきいん}} — \emph{ekiin} — \textquotedblleft{}a station attendant\textquotedblright{}.

\begin{quote}
\textbf{\ja{えきいんが} \ja{まどぐちに} \ja{います。}} — \emph{ekiin ga madoguchi ni imasu} — \textquotedblleft{}A station attendant is at the ticket window.\textquotedblright{}
\end{quote}

\begin{cousinweb}
\textbf{\ja{えきいん}} (\emph{ekiin}) is Sino-Japanese (kango): two Chinese characters read with their on'yomi, readings borrowed from Chinese. They are \emph{eki}, \textquotedblleft{}relay station\textquotedblright{} (Middle Chinese \emph{yek}) and \emph{in}, \textquotedblleft{}member\textquotedblright{} (\emph{hjwen}) (KANJIDIC2; Middle Chinese in Baxter's transcription).
\end{cousinweb}

\begin{grammarlens}[title={What you've built}]
One more word for everyday life at A2.
\end{grammarlens}

\subsection*{Your turn}
\begin{itemize}
\raggedright
  \item \emph{Say it:} \emph{ekiin}
  \item \emph{Say it:} \emph{ekiin}, once more
  \item \emph{Say it:} \emph{ekiin ga madoguchi ni imasu} — A station attendant is at the ticket window.
  \item \emph{Recall:} say the Japanese for lodging, then the Japanese for a ticket window, then say \emph{ekiin} again
\end{itemize}

\begin{tcolorbox}[breakable,colback=teal!4,colframe=teal!35!black,title={Before you move on}]
What does \ja{えきいん} mean? (\textquotedblleft{}A station attendant\textquotedblright{}.) Say it once more.
\end{tcolorbox}
\section[\texorpdfstring{\ja{かいさつぐち} (kaisatsuguchi)}{かいさつぐち (kaisatsuguchi)}]{\ja{かいさつぐち} (kaisatsuguchi) — a ticket gate}

\label{lesson:JA-C205-kaisatsuguchi}



\begin{quote}
Before the new one: say the Japanese for a ticket window, then the Japanese for a station attendant.
\end{quote}

\subsection*{You'll want to know: \ja{かいさつぐち}}
\textbf{\ja{かいさつぐち}} — \emph{kaisatsuguchi} — \textquotedblleft{}a ticket gate\textquotedblright{}.

\begin{quote}
\textbf{\ja{かいさつぐちは} \ja{ひだりです。}} — \emph{kaisatsuguchi wa hidari desu} — \textquotedblleft{}The ticket gate is on the left.\textquotedblright{}
\end{quote}

\begin{cousinweb}
\textbf{\ja{かいさつぐち}} (\emph{kaisatsuguchi}) is built from Sino-Japanese \emph{kaisatsu}, \textquotedblleft{}checking tickets\textquotedblright{}, read with on'yomi borrowed from Chinese, and native \emph{kuchi}, \textquotedblleft{}opening\textquotedblright{}, voiced to \emph{guchi} (JMdict, KANJIDIC2). The Chinese parts are \emph{kai}, \textquotedblleft{}to change, check\textquotedblright{} (Middle Chinese \emph{kojX}), and \emph{satsu}, \textquotedblleft{}tag, tablet\textquotedblright{} (\emph{tsreat}) (Baxter's transcription).
\end{cousinweb}

\begin{grammarlens}[title={What you've built}]
One more word for everyday life at A2.
\end{grammarlens}

\subsection*{Your turn}
\begin{itemize}
\raggedright
  \item \emph{Say it:} \emph{kaisatsuguchi}
  \item \emph{Say it:} \emph{kaisatsuguchi}, once more
  \item \emph{Say it:} \emph{kaisatsuguchi wa hidari desu} — The ticket gate is on the left.
  \item \emph{Recall:} say the Japanese for a ticket window, then the Japanese for a station attendant, then say \emph{kaisatsuguchi} again
\end{itemize}

\begin{tcolorbox}[breakable,colback=teal!4,colframe=teal!35!black,title={Before you move on}]
What does \ja{かいさつぐち} mean? (\textquotedblleft{}A ticket gate\textquotedblright{}.) Say it once more.
\end{tcolorbox}
\section[\texorpdfstring{\ja{かたみち} (katamichi)}{かたみち (katamichi)}]{\ja{かたみち} (katamichi) — one way; a single ticket}

\label{lesson:JA-C205-katamichi}



\begin{quote}
Before the new one: say the Japanese for a station attendant, then the Japanese for a ticket gate.
\end{quote}

\subsection*{You'll want to know: \ja{かたみち}}
\textbf{\ja{かたみち}} — \emph{katamichi} — \textquotedblleft{}one way; a single ticket\textquotedblright{}.

\begin{quote}
\textbf{\ja{かたみちの} \ja{きっぷを} \ja{ください。}} — \emph{katamichi no kippu o kudasai} — \textquotedblleft{}A one-way ticket, please.\textquotedblright{}
\end{quote}

\begin{cousinweb}
\textbf{\ja{かたみち}} (\emph{katamichi}) is native Japanese (wago), built from \emph{kata}, \textquotedblleft{}one side, one of a pair\textquotedblright{}, and \emph{michi}, \textquotedblleft{}road\textquotedblright{}, both inherited rather than borrowed (JMdict, KANJIDIC2).
\end{cousinweb}

\begin{grammarlens}[title={What you've built}]
One more word for everyday life at A2.
\end{grammarlens}

\subsection*{Your turn}
\begin{itemize}
\raggedright
  \item \emph{Say it:} \emph{katamichi}
  \item \emph{Say it:} \emph{katamichi}, once more
  \item \emph{Say it:} \emph{katamichi no kippu o kudasai} — A one-way ticket, please.
  \item \emph{Recall:} say the Japanese for a station attendant, then the Japanese for a ticket gate, then say \emph{katamichi} again
\end{itemize}

\begin{tcolorbox}[breakable,colback=teal!4,colframe=teal!35!black,title={Before you move on}]
What does \ja{かたみち} mean? (\textquotedblleft{}One way; a single ticket\textquotedblright{}.) Say it once more.
\end{tcolorbox}
\section[\texorpdfstring{\ja{おうふく} (ōfuku)}{おうふく (ōfuku)}]{\ja{おうふく} (ōfuku) — a round trip; a return ticket}

\label{lesson:JA-C205-oufuku}



\begin{quote}
Before the new one: say the Japanese for a ticket gate, then the Japanese for one way.
\end{quote}

\subsection*{You'll want to know: \ja{おうふく}}
\textbf{\ja{おうふく}} — \emph{ōfuku} — \textquotedblleft{}a round trip; a return ticket\textquotedblright{}.

\begin{quote}
\textbf{\ja{おうふくの} \ja{きっぷは} \ja{やすいです。}} — \emph{ōfuku no kippu wa yasui desu} — \textquotedblleft{}A return ticket is cheap.\textquotedblright{}
\end{quote}

\begin{cousinweb}
\textbf{\ja{おうふく}} (\emph{ōfuku}) is Sino-Japanese (kango): two Chinese characters read with their on'yomi, readings borrowed from Chinese. They are \emph{ō}, \textquotedblleft{}to go\textquotedblright{} (Middle Chinese \emph{hjwangX}) and \emph{fuku}, \textquotedblleft{}to return\textquotedblright{} (\emph{bjuwk}) (KANJIDIC2; Middle Chinese in Baxter's transcription).
\end{cousinweb}

\begin{grammarlens}[title={What you've built}]
One more word for everyday life at A2.
\end{grammarlens}

\subsection*{Your turn}
\begin{itemize}
\raggedright
  \item \emph{Say it:} \emph{ōfuku}
  \item \emph{Say it:} \emph{ōfuku}, once more
  \item \emph{Say it:} \emph{ōfuku no kippu wa yasui desu} — A return ticket is cheap.
  \item \emph{Recall:} say the Japanese for a ticket gate, then the Japanese for one way, then say \emph{ōfuku} again
\end{itemize}

\begin{tcolorbox}[breakable,colback=teal!4,colframe=teal!35!black,title={Before you move on}]
What does \ja{おうふく} mean? (\textquotedblleft{}A round trip; a return ticket\textquotedblright{}.) Say it once more.
\end{tcolorbox}
